\let\cleardoublepage\clearpage
\chapter*{Preface to the Fable Edition}
\addcontentsline{toc}{chapter}{Preface to the Fable Edition} % Add the new preface to the table of contents as a chapter

When I finished this book in the spring of 2023, ChatGPT was a few months old, ``prompt engineer'' sounded like a real job, and very serious people were signing very serious letters demanding that we pause AI research before it ended civilization. Three years later, civilization is still here, the letters have been forgotten, and the models have quietly become coworkers. It seemed like a good time to reopen the book and see how it held up.

This update is called the \textit{Fable Edition} because it was produced the way the book says work should now be produced: by a human expert and an AI working together. I gave an AI agent (built on a model named Fable) my manuscript, my bibliography, and my opinions, and it gave me back drafts, corrections, and the occasional argument. I kept final say on every word --- consider me the reasonably smart, not-too-cocky human getting recommendations from a model, exactly as Chapter 4 prescribes. If that arrangement sounds strange to you, this book is going to sound stranger.

A note on method, because honesty about method is half of this book's point. I have resisted the urge to silently rewrite my 2023 self. The original text is preserved wherever it held up, lightly corrected where it didn't, and the hindsight lives in the margins: wherever you see a blue \textcolor{fableblue}{\textbf{2026 update:}} note, that is the Fable Edition talking. Where the world changed enough to deserve more than a margin note --- self-driving cars, autonomous weapons, and the strange economics of open-source models --- I have added new sections, clearly dated, so you can tell the predictions from the postmortems.

The scoreboard, briefly. The boring claims aged best: that deep learning is a giant unscientific regression, that the data janitors matter more than the architects, that workers would adopt AI with or without permission, that open-source models would commoditize the technology and shrink its margins. The claim that aged worst is the one I was proudest of being skeptical about: I devoted a chapter to explaining why semi-autonomous driving is stupid, and then, somewhere between then and now, the statistics started winning. I bought a Tesla. Chapter 6 now ends with my confession, and I encourage you to enjoy it at my expense.

One thing has not changed, and it is the thesis of the book: the model proposes, the human disposes. The tools are shockingly better than they were in 2023 --- they write working code, they use other software, they run for hours without supervision --- and they are still complicated averages of their training data, still unable to explain themselves, and still most dangerous exactly where this book said they were: when we let them make critical decisions alone, and when we stop paying attention. The technology changed; the sucker traps didn't.

The professionally narrated audiobook still follows the first edition text, which I consider a feature: it is a time capsule of what we knew in 2023. This PDF is the living document. If you find something in here that has aged badly by the time you read it, write to me --- or have your agent write to mine.

\textit{Philadelphia, PA} \hfill \textit{Brad Flaugher}

\textit{June 2026}
